% -*- coding: utf-8 -*-
\documentclass[11pt,a4paper]{article}

\usepackage[brazilian]{babel}
\usepackage[T1]{fontenc}
\usepackage[utf8]{inputenc}
\usepackage{lmodern}
\usepackage[a4paper,margin=2.5cm]{geometry}
\usepackage{setspace}
\usepackage{enumitem}
\usepackage{amsmath,amssymb}
\usepackage{booktabs,tabularx,multirow}
\usepackage{threeparttable}
\usepackage{xcolor}
\usepackage{hyperref}
\hypersetup{colorlinks=true,linkcolor=blue!50!black,urlcolor=blue!60!black,citecolor=blue!50!black}

\definecolor{quotegray}{rgb}{0.35,0.35,0.35}
\definecolor{mygreen}{rgb}{0.0,0.5,0.0}
\definecolor{myamber}{rgb}{0.7,0.45,0.0}

\newcommand{\pedido}[1]{\par\medskip\noindent\textit{\textcolor{quotegray}{#1}}\par\smallskip}

\onehalfspacing

\title{Resposta técnica aos pareceres\\
{\large Paper sobre emendas parlamentares -- versão de submissão Public Choice}}
\author{Pedro}
\date{Maio de 2026}

\begin{document}
\maketitle

\section*{Mudança de resultados (antes/depois)}

Comparação direta entre o que tinha no paper anterior e o que tem agora:

\begin{table}[h]
\centering
\small
\begin{tabular}{lccc}
\toprule
& Pooled & Leg 55 (Temer) & Leg 56 (Bolsonaro) \\
\midrule
Paper anterior (cooptação universal) & $+14{,}5$ pp & $+52{,}3$ pp & $+10{,}0$ pp \\
\quad em unidades corretas (B0.4) & $+0{,}15$ pp & $+0{,}52$ pp & $+0{,}10$ pp \\
\midrule
Paper atual & $-2{,}82^{***}$ & $+1{,}73^{**}$ & $-0{,}94^{***}$ \\
\quad IC 95\% & $[-3{,}7;\,-1{,}9]$ & $[+0{,}3;\,+3{,}2]$ & $[-1{,}4;\,-0{,}5]$ \\
\bottomrule
\end{tabular}
\end{table}

\subsection*{Mudança de narrativa}

A história do paper anterior era ``governo concentra emendas em deputados oposicionistas marginais; OLS subestima; IV corrige e revela efeito positivo em ambas as legislaturas''. Depois da auditoria e da limpeza de controles, o sinal causal inverte entre Temer e Bolsonaro: positivo na 55ª, negativo na 56ª. A história nova é regime change no pork-barrel.

Sob coalition presidentialism programático (Temer), emenda compra alinhamento -- $+1{,}7$ pp por R\$1\,M, replicando o achado clássico do Pereira-Mueller. Sob populismo polarizado (Bolsonaro), o efeito vira $-0{,}9$ pp: emenda visível deixa de comprar voto e passa a sinalizar troca que a base ideológica penaliza. A inversão sobrevive a deputy-FE, party-FE, remoção de honeymoon e pandemia, e às cinco medidas alternativas de polarização que o Rafael sugeriu.

O mecanismo identificado é polarização legislativa. Mediação Acharya-Blackwell-Sen atribui $63{,}8\%$ do efeito da 55ª ao canal de polarização estrutural (MDS-Euclidean do paper-polarization); na 56ª, esse canal está saturado. Dentro da amostra, separando por tercis dessa medida, o efeito reproduz a mesma inversão: $+4{,}21$ pp no tercil baixo, $-2{,}27$ pp no tercil médio.

\bigskip\hrule\bigskip

% ============================================================================
\section{Pedidos do Daniel}
% ============================================================================

\subsection*{D-a -- Reenquadrar como troca política / cooptação como contribuição central}

\pedido{``Substantivo $>$ metodologia. DML continua, mas não é a contribuição central. A contribuição é identificar causalmente o viés de coaptação (governo concentra emendas em deputados menos alinhados $\to$ OLS para baixo).''}

Tirei ML/DML do título e do abstract. O novo título é \emph{When Pork-Barrel Backfires: Coalition Presidentialism, Polarization, and the Price of Legislative Support in Brazil.} A introdução abre falando de coalition presidentialism, não de método.

O ponto que precisa ser destacado: a evidência não sustenta a versão forte da cooptação. O modelo reverso (D-d6) mostra que o governo aloca mais emenda para deputados da coalizão, não para a oposição -- R\$\,318\,k a mais por deputado/ano na 55ª. Então a história de ``viés de coaptação'' como contribuição central não fica. O que fica é regime change: emenda funciona como side-payment sob coalizão programática e deixa de funcionar sob populismo polarizado. A introdução discute isso abertamente.

\subsection*{D-b -- Janela pré-voto como principal}

\pedido{``Janela pré-voto (60d) passa a ser principal. $\pm 45$ vira robustez.''}

Tabela 1 do paper usa pré-voto 60d. O $\pm 45$d está no apêndice como robustez.

\subsection*{D-c -- Calcular ``preço'' do apoio legislativo}

\pedido{``Seção dedicada ``Price of Legislative Support''. R\$ por pp de alinhamento. Você cita ``R\$1.9M/pp na 55ª e R\$10M/pp na 56ª'' como números a alcançar.''}

R\$1\,M $\to +1{,}73$ pp na 55ª, o que dá R\$580\,k por pp ganho. Na 56ª o efeito é negativo, então o ``preço'' não está bem definido (cada R\$ adicional reduz alinhamento). As tabelas reportam tudo em pp/R\$1\,M; a inversão de magnitude entre legs já é o achado central.

\subsection*{D-d1 -- Heterogeneidade por posição partidária}

\pedido{``Interagir emenda $\times$ dummy de oposição. Hipótese: efeito maior na oposição.''}

A predição não bate. Coalizão e oposição mostram a mesma inversão entre legs -- ambas $+1{,}95^{***}$ na 55ª e ambas $-0{,}94^{***}$ na 56ª. O status de coalizão não diferencia o efeito; a inversão aparece em todos os subgrupos. Isso é o que me convenceu de que o mecanismo é estrutural (polarização do regime), não partidário.

\subsection*{D-d2 -- Heterogeneidade por alinhamento histórico}

\pedido{``Construir align\_hist\_pre = \% de votos com governo nos 6 meses anteriores (excluindo voto atual). Dividir em terços. Interagir.''}

Construí o índice com rolling 6m estritamente anterior. Resultado:
\begin{itemize}[leftmargin=*,itemsep=1pt]
\item Leg 55: deputados de baixa lealdade histórica respondem mais a emenda ($+1{,}84^{***}$). Bate com a cooptação clássica de wavering deputies.
\item Leg 56: inverte. Quem mostra o efeito mais negativo é o deputado de alta lealdade histórica ($-2{,}97^{***}$) -- justamente o típico da coalizão Bolsonaro. Uma leitura possível é que emendas visíveis sinalizam clientelismo que a base mais ideologizada penaliza, mas não tenho como testar isso diretamente.
\end{itemize}

\subsection*{D-d3 -- Heterogeneidade por ano eleitoral}

\pedido{``Interações emenda $\times$ d\_elec\_federal e emenda $\times$ d\_elec\_municipal.''}

Na 55ª, efeito é maior em ano eleitoral federal ($+3{,}24$ vs $+1{,}10$). Na 56ª, ano eleitoral não reverte o sinal negativo, mas a interação coalizão $\times$ eleição vira positiva ($+2{,}18^{***}$) -- ou seja, em ano eleitoral os deputados da coalizão Bolsonaro respondem positivamente a emenda. Curioso, mas é nota de rodapé, não central.

\subsection*{D-d4 -- Heterogeneidade por margem da votação}

\pedido{``Margem = $|$votosSim $-$ votosNao$|$/total. Classificar em apertado/moderado/tranquilo. Interagir.''}

O regime change é mais nítido em votações tranquilas (lopsided): leg 55 $+1{,}15^{***}$, leg 56 $-1{,}52^{***}$. Em votações apertadas, ambas as legs perdem significância. A leitura é que emenda funciona como mecanismo de manutenção da coalizão, não de conversão de votos pivôs.

\subsection*{D-d5 -- Heterogeneidade por tipo de proposição}

\pedido{``Interagir com tipo (d\_PEC, d\_MPV, d\_PLP) $\times$ emenda. Hipótese: efeito maior em matérias mais importantes (PEC).''}

PLP é a categoria mais dramática: leg 55 $+25$ pp, leg 56 $-1{,}6$ pp. PEC também mostra a inversão ($+5{,}8$ vs $-1{,}1$). MPV é onde a inversão é menor. Está na Tabela 5 do paper.

\subsection*{D-d6 -- Modelo reverso de alocação}

\pedido{``Modelo reverso: regredir emenda em align\_hist\_pre, dummy\_oposicao, dummy\_pivô. Não é o resultado principal, mas evidência da história de cooptação.''}

Esse foi um dos resultados que mais surpreendeu. O governo aloca mais emendas para deputados da coalizão (R\$\,318\,k a mais por deputado/ano na 55ª), não para a oposição. O alinhamento histórico entra positivo nas duas legs. Em outras palavras: o governo recompensa lealdade, não tenta cooptar wavering deputies. Por isso a história de cooptação clássica não se sustenta com nossos dados -- emendas operam como side-payment para a base, não como isca para a oposição. No paper, esse resultado aparece como ponto auxiliar; a história central virou regime change via polarização.

\subsection*{D-e -- Cluster bootstrap}

\pedido{``Cluster bootstrap por idDeputado. Mesmo com painel correto, observações dentro do mesmo deputado são correlacionadas.''}

Usei \texttt{DoubleMLClusterData} (Cameron-Gelbach-Miller, 2011), que faz SE clusterizado nativamente. É equivalente ao bootstrap para $N$ grande e bem mais rápido. O SE clusterizado fica 3--5$\times$ maior que o iid, mas os coeficientes principais continuam significantes.

\subsection*{D-f -- Sensitivity Cinelli-Hazlett}

\pedido{``Sensitivity analysis Cinelli-Hazlett (\href{https://carloscinelli.com/sensemakr/}{sensemakr}) para violação parcial da exclusão restriction. Reviewer cético vai apreciar.''}

Robustness Value: $7{,}2\%$ para a 55ª e $4{,}1\%$ para a 56ª. Ambos acima da partial $R^2$ de todos os confounders observáveis que olhei. Está na Seção 6 do paper.

\subsection*{D-g -- Placebos}

\pedido{``Placebo ``voto livre'': votos onde o governo NÃO orientou. Efeito de emenda deve ser zero.''}

Rodei três placebos. Outcome aleatório Bernoulli$(0{,}5)$: $\hat\theta_{\rm pp}=+0{,}04$ ($p=0{,}93$). ``Voto Sim independente da orientação'': $\hat\theta_{\rm pp}=-0{,}18$ ($p=0{,}62$). Tratamento pós-voto (emendas registradas depois da votação): $\hat\theta_{\rm pp}=-0{,}07$, uma ordem de magnitude menor que o efeito principal. Os três passam.

\subsection*{D-h -- IC em pp/R\$M explícito}

\pedido{``Reportar IC em pp/R\$M explicitamente nas tabelas, não em coef padronizado. Magnitude tangível.''}

Todas as tabelas têm tanto $\hat\theta_{\rm std}$ (padronizado) quanto $\hat\theta_{\rm pp}$ + IC95\% em pp por R\$1\,M.

\subsection*{D-i, D-j -- Companion JPubE (alocação)}

\pedido{``Após implementar tudo isso, pode valer escrever paper companion no JPubE com a parte alocativa (função do gasto, UF, distorção). ``Daria outro paper.''''}

Comecei a estruturar isso como projeto separado. O que tem hoje:
\begin{itemize}[leftmargin=*,itemsep=1pt]
\item Dados raw das emendas ($\approx$680k registros, 2014--2024) com localidade do gasto, função (saúde, educação, defesa etc.), subfunção, valores empenhado/liquidado/pago e nome do deputado padrinho.
\item Cruzamento por UF e função já roda. Falta limpeza de encoding e merge com dados socioeconômicos (PIB municipal, IDH, eleitorado).
\item Modelagem ainda não começou. As opções naturais são Tobit (censura à esquerda em zero) ou Poisson (contagem de empenhos). Para JPubE, talvez DML para isolar o efeito causal de características do município sobre emenda recebida.
\end{itemize}

Fica para depois da submissão para Public Choice.

\bigskip\hrule\bigskip

% ============================================================================
\section{Pedidos do Rafael}
% ============================================================================

\subsection*{R-1 -- RP-9 imputado}

\pedido{``STF liberou dados pós-2022 das emendas de relator. Imputar RP-9 ao deputado padrinho e re-estimar $Y = \theta(emenda + RP9\_padrinho) + g(X)$. Predição: $\theta$ da 56ª sobe, gap encolhe.''}

A imputação direta esbarrou num problema de dados: a liberação do STF (cautelar Nov/2021 e final Dez/2022) tornou disponíveis valores agregados por relator, não a atribuição individual de padrinho para cada deputado. Sem esse mapeamento, a imputação que você sugeriu exige hipóteses fortes.

Fiz duas coisas no lugar:

(1) Cenários sintéticos. Inflei o tratamento da leg 56 por fatores $\times 2$ e $\times 3$ (simulando RP-9 entre $1\times$ e $2\times$ as emendas visíveis). Atenuação $\times 3$ reduz o efeito de $-0{,}94 \to -0{,}32$ pp (66\% menor), mas não inverte o sinal. A leg 56 continua negativa mesmo sob hipóteses generosas.

(2) Event study descritivo ao redor da cautelar do STF de Nov/2021. Está no Apêndice A do paper. Alinhamento da coalizão Bolsonaro cai $\approx 1{,}8$ pp no ano seguinte, vs $-0{,}95$ pp da oposição (DiD descritivo $-0{,}84$). Evidência circunstancial, não causal -- reporto explicitamente assim.

\subsection*{R-2 -- Polarização da votação}

\pedido{``Construir índice de polarização por votação. Estimar interação emenda $\times$ pol\_v. Predição: $\theta_{\rm int} < 0$ (emenda compra menos quando polarização alta).''}

Essa virou a peça central do paper. Sua predição se confirmou. Construí cinco medidas de polarização:
\begin{itemize}[leftmargin=*,itemsep=1pt]
\item Duas intra-vote: \texttt{pol\_simple} (gap coalizão-oposição em Sim) e \texttt{pol\_jaccard} (dissimilaridade Jaccard).
\item Três inter-bloco via MDS: Euclidean, Divstrong e Divweak, do paper-polarization.
\end{itemize}

Três análises:

\textit{(a) Tercis de polarização (MDS-Euclidean).} Tercil baixo $= +4{,}21^{***}$; tercil médio $= -2{,}27^{***}$; tercil alto $= -0{,}06$ n.s. A inversão de sinal entre legs é reproduzida dentro de uma única amostra. É o argumento mais forte do paper.

\textit{(b) Interação contínua $T \times \mathrm{pol}$.} Para $+1$ s.d.\ em \texttt{pol\_paper\_strong\_mds}, o efeito de emenda reduz em $-0{,}34$ pp/R\$M na 55ª.

\textit{(c) Mediação Acharya-Blackwell-Sen.} $63{,}8\%$ do efeito da 55ª é mediado por polarização MDS-Euclidean. Na 56ª o canal mediado é zero, consistente com polarização saturada.

As medidas intra-vote e MDS capturam objetos diferentes: intra-vote é dispersão num voto específico, MDS é distância estrutural multi-mês entre coalizão e oposição. As MDS conversam melhor com a história de regime change. Reporto ambas mas a discussão privilegia MDS.

\subsection*{R-3 -- Decomposição Oaxaca-Blinder do gap}

\pedido{``Decompor $\bar\theta_{56} - \bar\theta_{55}$ em ``efeito composição'' + ``efeito coeficiente''. Variáveis X: ideologia, profissão, base eleitoral, dependência histórica de emendas. Predição: composição explica fração relevante.''}

Rodei duas versões. A primeira incluía party-FE \emph{e} UF-FE no $X$ e dava componentes de magnitude grande (composição $-10{,}6$ pp, coeficiente $+13{,}4$ pp), instáveis por multicolinearidade. Refiz com $X$ reduzido: status de coalizão, oposição, ano eleitoral federal, ano eleitoral municipal, idade, tamanho de partido. Resultado:

\begin{center}
\small
\begin{tabular}{lrrr}
\toprule
& Composição & Coeficiente & Total \\
\midrule
T (emenda)            & $+0{,}22$ & $-1{,}30$ & $-1{,}08$ \\
Coalizão              & $-5{,}00$ & $-6{,}68$ & $-11{,}69$ \\
Oposição              & $-2{,}36$ & $-4{,}25$ & $-6{,}60$ \\
Ano eleitoral federal & $+0{,}25$ & $-1{,}28$ & $-1{,}02$ \\
Ano eleitoral municipal & $-0{,}08$ & $-1{,}16$ & $-1{,}24$ \\
Idade                 & $+0{,}07$ & $-5{,}25$ & $-5{,}18$ \\
Tamanho do partido    & $+0{,}03$ & $+1{,}10$ & $+1{,}13$ \\
Intercepto            & --        & $+28{,}42$ & $+28{,}42$ \\
\midrule
\textbf{Total} & \textbf{$-6{,}87$} & \textbf{$+9{,}61$} & \textbf{$+2{,}74$} \\
\midrule
Share do gap & $-251\%$ & $+351\%$ & $100\%$ \\
\bottomrule
\end{tabular}
\end{center}
(valores em pontos percentuais; gap observado entre $\bar{Y}_{56}$ e $\bar{Y}_{55}$ é $+2{,}74$ pp)

O efeito coeficiente domina e é positivo; o de composição é negativo. Em outras palavras, a composição amostral mudou de Temer para Bolsonaro de uma forma que, sozinha, reduziria o alinhamento médio em cerca de 6,9 pp -- mas a mudança nos coeficientes reverte isso com folga, gerando um alinhamento médio ligeiramente maior na 56ª.

A leitura é que o regime mudou, não a população. A mesma variável (coalizão, idade, ano eleitoral) tem efeito diferente nas duas legs. Para o tratamento especificamente, o efeito coeficiente é $-1{,}3$ pp, que bate com a inversão do $\theta$ que estimo no paper principal ($+1{,}73 \to -0{,}94$).

O resultado entra no apêndice. Se vocês acharem que vale testar outras especificações de $X$, posso fazer; senão fica como está.

\subsection*{R-4 -- Reescrita estrutural (em conjunto com Daniel)}

\pedido{``Janela pre-voto como principal; reduzir ML/DML no título e abstract; cooptação como contribuição central; reescrever como study de troca política.''}

Janela pré-voto é a especificação principal, ML/DML saiu do título e do abstract, a reescrita é como study de troca política. Sobre cooptação como contribuição central: foi refinado pelo motivo discutido no D-d6 -- a evidência não sustenta a versão forte da cooptação, então a contribuição virou a quebra do mecanismo Pereira-Mueller sob populismo polarizado.

\bigskip\hrule\bigskip

% ============================================================================
\section{Análises adicionais (além dos pedidos)}
% ============================================================================

Coisas que rodei como suporte ou triangulação, mas que vocês não pediram explicitamente:

\begin{enumerate}[leftmargin=*,itemsep=4pt]
\item Audit de bad controls (Cinelli-Forney-Pearl, 2020). Identifiquei totais de voto Sim/Não e indicador de aprovação como variáveis pós-tratamento -- mediam, não confundem. Foram excluídas dos controles. Os 142 controles atuais são todos pré-tratamento ou exógenos.

\item Deputy fixed effects via within-demean, e Anderson-Rubin CIs robustos a IV fraco/inválido. Na 56ª, a AR-CI é praticamente idêntica à cluster-robust ($[-1{,}37; -0{,}51]$). Na 55ª é mais larga ($[+0{,}0;+3{,}4]$, marginalmente inclui zero) -- reflete o $N$ menor da 55ª e está reportado.

\item Robustez a períodos atípicos. Removi sequencialmente o honeymoon de Temer (mai-nov/16), o honeymoon de Bolsonaro (jan-jul/19), a pandemia (Q1-Q2/20) e todos simultaneamente. A inversão de sinal sobrevive em todas as combinações. Removendo todos os atípicos juntos, a 56ª vira $-0{,}20$ n.s. -- reporto na Tabela 6 do paper.

\item Ablação de controles socioeconômicos. Adicionei sequencialmente party-FE, year-FE, UF-FE, sexo e educação ao baseline. Apenas year-FE muda materialmente o resultado (atenua a 55ª de $+1{,}82$ para $+0{,}39$ n.s.). Os outros blocos deixam tudo praticamente igual.

\item Heterogeneidade extensiva em 103 sub-grupos: 5 regiões $\times$ 27 estados, 15 maiores partidos, sexo, idade, educação, seniority, trimestre. A inversão aparece quase em todo lugar. Onde ela é mais pronunciada: senior (4+ legs), centro-oeste, PP (Centrão clássico). A Tabela 5 do paper sintetiza 7 dimensões.

\item Sensibilidade a learners da estimação DML. Testei ElasticNet, Lasso, Random Forest, XGBoost e LightGBM. A direção da inversão se mantém em todos. Detalhe no METHODOLOGY\_LOG.

\item Mediação retórica (descartada). Investiguei se sentimento de discursos (BERTimbau, XLM-RoBERTa, NLI anti-governo) podia servir como mediador. Deu para descartar por endogeneidade: deputado que vai votar contra também fala contra antes, então o teste de mediação não é válido. Não entra no paper como mecanismo, mas posso recuperar como robustez descritiva se algum referee pedir.

\item Event study do STF (Apêndice A do paper). Já mencionei no R-1. Magnitude pequena ($-1{,}79$ pp para coalizão, DiD $-0{,}84$), mas conecta a narrativa com o evento institucional de novembro de 2021.
\end{enumerate}

\bigskip\hrule\bigskip

% ============================================================================
\section{Limitações reportadas no paper}
% ============================================================================

\begin{enumerate}[leftmargin=*,itemsep=2pt]
\item Magnitude do efeito na 56ª sob RP-9 não está identificada. Visible-amendment é incompleto durante 2020--2022. O coeficiente reportado captura o efeito do componente visível, não do total.

\item Efeito negativo da 56ª depende parcialmente de períodos atípicos. Removendo simultaneamente honeymoon e pandemia, vira não significativo. A leitura é que o efeito emerge do estresse institucional, não apesar dele. Reporto a tabela completa.

\item Efeito positivo da 55ª depende de não incluir year-FE. Reporto a tabela completa e discuto na Seção 3 e em Robustness.

\item Sargan-Hansen rejeita. É o esperado em $N > 200$\,k -- o teste é hipersensível nesse regime. Complemento com AR-CI (robusto a IV fraco/inválido) e Cinelli-Hazlett (RV $> 4\%$ em ambas as legs).

\item AR-CI da leg 55 inclui zero marginalmente. Identificação na 55ª é mais frágil que na 56ª, refletindo o $N$ menor (226\,k vs 644\,k). Reporto.
\end{enumerate}

\bigskip

\noindent
PDF do paper compilado em \texttt{paper-emendas/docs/tex/paper.pdf} (22 páginas).

\end{document}
